\cotefile{27}
\scanpage{27}
\deleted{\uline{Corollaire} 1 \quad Soit $\unsure{S|g}$ une section de $G$}

\noindent\begin{tabular}{@{}|p{.94\linewidth}@{}}
\deleted{(i) Si \deleted{\unsure{$g$ est}} $s\in S$ est tel que $g(s)$ \unsure{régulier|regulier}, alors pour $s'$ voisin de $s$, $g(s')$ est régulier.}
\end{tabular}

\noindent\begin{tabular}{@{}|p{.94\linewidth}@{}}
\deleted{(ii) Si $g$ est régulier, alors $\exists$ un sous-groupe de Cartan et un seul $C$ de $G$ qui contient $g$.}
\end{tabular}

\deleted{\uline{\unsure{Corollaire|Contrexemple}} 2}
% archive sheet number: 26
\scanpage{28}
\uline{Théorème de $1/2$ continuité}
% archive sheet number: 27
\scanpage{29}
$G$ lisse connexe sur $k$ alg. clos.

\diagram{\begin{tikzcd}[ampersand replacement=\&,column sep=small,row sep=0pt]
e \arrow[r,dash] \& H \arrow[r,dash] \& G \arrow[r] \& A \arrow[r,dash] \& e \\
\& \substack{\text{groupe affine}\\\text{lisse conn.}} \& \& \text{abélienne} \&
\end{tikzcd}}

\[
\begin{aligned}
\rho_{ab}(G)&=\dim A & \rho_{aff}(G)&=\dim H\\
\rho_r(G)&=\rho_r(H)=\dim\text{ tores maximaux}
\end{aligned}
\]
\margin{\(\left\{\begin{aligned}
\rho_{\deleted{\unsure{s}}rs}(G)&=\rho_{rs}(H)=\\[-2pt]
&\rho_r(H/\mathrm{rad}(H))\\
\deleted{\rho_{rad}(G)=} &\deleted{\rho_{rad}(H)=\unsure{\rho_r}}\\
&\deleted{\rho_r(\mathrm{rad}(H))}\\
\rho_{rad\deleted{\unsure{r}}}(G)&=\rho_r
\end{aligned}\right.\)}

\[
\begin{aligned}
\rho_s(G)&=\rho_s(H)=\deleted{\unsure{\rho_r}}\ \deleted{\dim}\rho_r(H/\mathrm{rad}(H))\\
d_s(G)&=d_s(H)=\dim(H/\mathrm{rad}(H))\\
\rho_n(G)&=\rho_n(V)+\rho_{ab}(G)=\text{dimension des Cartans}\\
\rho_{na}(G)&=\rho_n(V)\\
\rho_u(G)&=\rho_u(V)=\rho_{na}(G)-\rho_r(H)
\end{aligned}
\]

\[
\deleted{\unsure{\rho_{aff}}},\ \underline{\rho_{ab}},\
\left\lfloor\underline{\rho_r,\ \deleted{\unsure{\rho_s}},\ \rho_{u\deleted{\unsure{a}}}}\right\rfloor
\left\lceil\overline{,\ \underline{\unsure{\phi_{ss}|\rho_{ss}|d_{ss}}},\ \underline{\unsure{\phi_{rad}|\rho_{rad}|d_{rad}}}}\right\rceil
\qquad
\deleted{\dim G=\rho_{aff}+\rho_{ab}=\underline{\rho_{rad}}+\rho_{ss}+\rho_{ab}}
\]

\[
\begin{aligned}
\rho_{na}&=\rho_r+\rho_u & \underline{\rho_{aff}},\ \rho_{rad}\\
\rho_n&=\rho_{na}+\rho_{ab}=\rho_r+\rho_u+\rho_{ab}
\end{aligned}
\]

\[
\begin{aligned}
\deleted{\rho_s\geq\rho_{rs}\unsure{=}}&\\
\deleted{\rho_r=\rho_{rs}+\rho_{rr}}&
\end{aligned}
\quad
\left\{\begin{aligned}
\rho_{aff}&=\rho_{ss}+\rho_{rad}\\
\dim G&=\rho_{ab}+\rho_{aff}=\rho_{ab}+\rho_{ss}+\rho_{rad}
\end{aligned}\right.
\]

\noindent\begin{tabular}{@{}|p{.94\linewidth}@{}}
$\rho_r,\rho_u,\rho_{ab},\rho_{na},\rho_n$ sont les $\overline{\mathrm m}\mathrm{s}$ pour $G$ et pour un ss-groupe \unsure{contenant} un Cartan \ldots
\end{tabular}
% archive sheet number: 28
\scanpage{30}
\[
\deleted{\unsure{n A_0|{}_n A_0}}
\]
\[
\begin{aligned}
\deleted{T_0\subset C_0} &\qquad \deleted{(x u)}\\
\deleted{T_0\subset C_0^{aff}\underbrace{\subset C_0}_{\mathrm{V.A}}} &\qquad \deleted{x^n=1}\\
&\qquad\deleted{(xu)^n={}?\,?}
\end{aligned}
\]
\[
\begin{aligned}
\deleted{H^2(\ ,T_0)\simeq H^2({}_{\infty}T_0)={}}\\
\deleted{H^2(\ ,\unsure{Q|\mathbb{Q}}/\mathbb{Z})}
\end{aligned}
\]
\[
\begin{aligned}
\deleted{G\quad -} &&& \deleted{\rho_{ab}\leq\rho'_{ab}}\\
\deleted{\mathrm{I}\quad\mathrm{I}} &&\deleted{\begin{smallmatrix}\rho_{ab}\ \deleted{\unsure{+\rho_r}}\\ \rotatebox{90}{$\leq$}\\ \rho'_{ab}\end{smallmatrix}}&
\deleted{\frac{G_1}{Z_1}=Z}\\
\deleted{\rho_{ab}+2\rho_r} &&&\deleted{\rho_{ab}(Z_0)\leq\rho_{ab}(Z_1)=\rho_{ab}(G_1)}
\end{aligned}
\]
\[
\deleted{\rho_a+\qquad\rho_{ab}+}
\]

\diagram{\begin{tikzcd}[ampersand replacement=\&,column sep=2pt,row sep=2pt,execute at end picture={
\draw (b.north west) ++(0,2pt) -- ++(0,4pt) -| ([yshift=2pt]c.north east);
\draw (b.north west) ++(-2pt,7pt) -- ++(0,10pt) -| ([yshift=6pt]e.north west);
\draw (e.north west) ++(2pt,2pt) -- ++(0,5pt) -| ([yshift=2pt]f.north east);
\draw ([yshift=7pt]e.north) -- ++(0,-4pt);
\draw ([yshift=-2pt]b.south west) -- ++(0,-4pt) -| ([yshift=-2pt]d.south east);
\draw ([yshift=-2pt]c.south) -- ++(0,-4pt);
\draw ([yshift=-2pt]d.south east) -- ++(2pt,-9pt) -| ([yshift=-2pt]f.south);
\draw ([yshift=-3pt]e.south) -- ++(2pt,-22pt) -| ([yshift=-3pt]f.south east);
}]
\& \& \rho_{\mathrm{rig}.} \& \& \& \\
\& \& \rho_{n\,aff.} \& \& d \& \\
\rho_{ss} \& |[alias=b]|\underset{\cdot}{\rho_{rr}} \& |[alias=c]|\underset{\cdot}{\rho_u} \& |[alias=d]|\underset{\cdot}{\rho_{ab}} \& |[alias=e]|\underset{\cdot}{d_{ss}} \& |[alias=f]|\underset{\cdot}{d_{rad\,aff}}\\
\& \& \rho_n \& \& d_{rad}\& \\
\& \& \& \& \& d_{aff}
\end{tikzcd}}

\[
\begin{array}{c}
\text{\textcircled{\scriptsize 1}}\\\text{\textcircled{\scriptsize\unsure{1|2}}}\\\text{\textcircled{\scriptsize 3}}\\\text{\textcircled{\scriptsize 4}}
\end{array}
\left\{\begin{array}{l}
\tikz[baseline=(v.base)]{\node[draw,rounded corners=8pt,inner sep=2pt](v){$\rho_{ab}$};}\\
\tikz[baseline=(v.base)]{\node[draw,rounded corners=8pt,inner sep=2pt](v){$\rho_r+\rho_{ab}=\rho_{\mathrm{rig}}$};}\\
\tikz[baseline=(v.base)]{\node[draw,rounded corners=8pt,inner sep=2pt](v){$\rho_{ss}$};}\\\tikz[baseline=(v.base)]{\node[draw,rounded corners=8pt,inner sep=2pt](v){$d_{ss}$};}
\end{array}\right.
\qquad
\left\{\begin{array}{l}
\tikz[baseline=(v.base)]{\node[draw,rounded corners=8pt,inner sep=2pt](v){$\rho_n$};}\\\rho_u\\\rho_{n\,aff.}\\d_{aff}\\d_{rad}\\d_{rad\,aff.}
\end{array}\right.
\begin{array}{c}
\text{\textcircled{\scriptsize 5}}\\\text{\textcircled{\scriptsize 6}}\\\text{\textcircled{\scriptsize 7}}\\\text{\textcircled{\scriptsize 8}}\\\text{\textcircled{\scriptsize 9}}\\\text{\textcircled{\scriptsize 10}}
\end{array}
\]

\[
\deleted{\rho_r=\rho_{ss}+\deleted{\unsure{\rho_{rad}}}\rho_{r.\,rad\,aff}}
\]
\[
d_{rad\,aff}\geq\rho_u+\widetilde{\rho}_r
\]
\[
\boxed{\rho_{ab}+\unsure{d_{ss}|\phi_{ss}}+\underbrace{(\rho_r-\rho_{ss})}_{\rho_{r\,rad\,aff}}}
\quad \rho_{ab}+\rho_r+(d_{ss}-\rho_{ss})
\]
\[
\begin{aligned}
\text{\textcircled{\scriptsize 2}}+\text{\textcircled{\scriptsize 5}}&\Longrightarrow\text{\textcircled{\scriptsize 6}}\\
\text{\textcircled{\scriptsize 1}}+\text{\textcircled{\scriptsize 5}}&\Longrightarrow\text{\textcircled{\scriptsize 7}}\\
\text{\textcircled{\scriptsize\unsure{1|4}}}&\Longleftrightarrow\text{\textcircled{\scriptsize 8}}\\
\text{\textcircled{\scriptsize 3}}&\longleftrightarrow\text{\textcircled{\scriptsize 9}}\\
\text{\textcircled{\scriptsize 1}}+\text{\textcircled{\scriptsize 3}}&\Longrightarrow\text{\textcircled{\scriptsize 10}}
\end{aligned}
\]
% archive sheet number: 29
\scanpage{31}
\mspage{\handwritten{\unsure{63}}}
signaler en remarque). Par suite, $\operatorname{Ker}i$ est isomorphe à $N_{k(\underline p)/k}$, d'où
\[
\operatorname{rang}\operatorname{Im}i=r-t=n-s-t
\]
où
\[
\begin{aligned}
n&=\dim A,\quad \handwritten{s}=\dim A/\underline p,\quad r=n-s=\dim A_{\underline p}=\operatorname{rang}\underline p/\underline p^{2}\otimes k(t)\\
t&=\operatorname{rang}N_{k(p)/k}
\end{aligned}
\]
D'autre part, la deuxième ligne de $(D_{\handwritten{\alpha}})$ nous prouve que
\[
\operatorname{rang}\operatorname{Im}i_{\handwritten{\alpha}}
=\operatorname{rang}\handwritten{\hat\Omega}^{1}_{A/K_{\handwritten{\alpha}}}\otimes k(p)
-\operatorname{rang}\handwritten{\hat\Omega}^{1}_{k(\underline p)/K_{\handwritten{\alpha}}}
\]
Or \deleted{le premier terme du second membre, compte tenu que $A$ est \unsure{formellem|formelle}} sauf erreur il doit être vrai que pour $\handwritten{\alpha}$ grand, $\handwritten{\hat\Omega}^{1}_{A/K_{\handwritten{\alpha}}}$ est un module \handwritten{\uline}{libre de rang} $n+\operatorname{rang}\handwritten{\Omega}^{1}_{K_{\unsure{o|0}}/K_{\handwritten{\alpha}}}$ sur $A$. Il est vrai que je ne l'ai pas prouvé, et \deleted{\unsure{étant}} voulant en finir, je vais faire des hypothèses plus restrictives dans l'énoncé de 8.1. savoir que $A$ est formellement simple, non seulement sur $k$, mais même sur $K_o$ (il faut donc modifier la rédaction de l'énoncé 8.1. à moins que tu ne te sentes le courage de résoudre le point que je laisse en suspens ici). Dans ce cas, l'assertion que j'ai faite sur $\handwritten{\hat\Omega}^{1}_{A/K}$ est même vrai pour \uline{tout} sous-corps $K$ d'indice fini dans $K_o$, et se vérifie sans difficulté. \deleted{\unsure{E|K}}\unsure{C|K}omparant alors les deux expressions trouvées pour $\operatorname{rang}\operatorname{Im}i_{\handwritten{\alpha}}$ et $\operatorname{rang}\operatorname{Im}i$, on est remené à montrer que les deux sont égaux pour $\handwritten{\alpha}$ grand, i.e. qu'on a pour $\handwritten{\alpha}$ grand l'égalité
\[
\operatorname{rang}\handwritten{\hat\Omega}^{1}_{k(\underline p)/K_{\handwritten{\alpha}}}
-\operatorname{rang}\handwritten{\Omega}^{1}_{K_o/K_{\handwritten{\alpha}}}
=s+t=\dim A/\underline p+\operatorname{rang}N_{\handwritten{k(p)/k}}
\]
Pour abréger, posons $L=k(\underline p)$, \handwritten{no}tons que c'est une extension finie d'une extension de la forme $\handwritten{\bar}{K}_o=K_o((y_1,\ldots,y_s))$, et que l'on a
\[
\handwritten{\hat\Omega}^{1}_{L/K_{\handwritten{\alpha}}}=\handwritten{\Omega}^{1}_{L/K'_{\handwritten{\alpha}}}
\]
où on pose
\[
K'_{\handwritten{\alpha}}=K_{\handwritten{\alpha}}((y_1^p,\ldots,y_s^p)),
\]
ce qui ramène les modules avec chapeau à des modules sans. On est ainsi
% archive sheet number: 30
\scanpage{32}
\uline{Th} Soit $G$ un préschéma en groupes \inserted{lisse} séparé de présentation finie sur $S$. Alors

\noindent\begin{tabular}{@{}|p{.94\linewidth}@{}}
(i) \deleted{les fonctions} \deleted{\unsure{Th|les}} les fonctions \inserted{$\rho_{aff},\deleted{\unsure{\rho}}d\deleted{\dim G}$ sont $1/2$ continues sup, et de} $\rho_n,\rho_{na},\rho_u$ sont semi-continues \uline{supérieurement} \deleted{si $G$ est plat sur $S$ [\unsure{alors|donc} $d$ est \unsure{lct.|loc.}\ \unsure{cte|c\textsuperscript{st}}]}.
\end{tabular}

\noindent\begin{tabular}{@{}|p{.94\linewidth}@{}}
(ii) \deleted{Si $G$ est \uline{plat} sur $S$,} \deleted{les} fonctions $\rho_{ab},\rho_{ab}+\rho_r,\rho_s$ sont semi-continues inférieurement.
\end{tabular}

On se ramène comme d'habitude au cas $S$ affine, \unsure{noth|noeth}, spectre d'un anneau de valuation discrète, pt fermé $s_0$ pt générique $s_1$, et à prouver les inégalités
\[
\left\{\begin{aligned}
&\inserted{\rho_{aff}\geq\rho'_{aff}\qquad d\geq d'}\\
&\deleted{\text{\textcircled{\scriptsize\unsure{4}}}}\text{\textcircled{\scriptsize 5}}\rho_{\deleted{\unsure{s}}n}\geq\rho'_n,\quad
\deleted{\text{\textcircled{\scriptsize\unsure{5}}}}\text{\textcircled{\scriptsize 6}}\rho_u\geq\rho'_u\quad
\text{\textcircled{\scriptsize 7}}\rho_{na}\geq\rho'_{na}\quad\deleted{,\rho_{aff}\geq\rho'_{aff},}\\
&\text{\textcircled{\scriptsize 1}}\rho_{ab}\leq\rho'_{ab},\quad \text{\textcircled{\scriptsize 2}}\rho_{ab}+\rho_r\leq\rho'_{ab}+\rho'_r,\quad \text{\textcircled{\scriptsize 3}}\rho_s\leq\rho'_s
\end{aligned}\right.
\]
\deleted{\unsure{les}}\unsure{les} \deleted{deux}\inserted{quatre} premières sans réserve de platitude.

Ce sont \deleted{\unsure{des}} (pour un groupe algébrique sur un corps alg clos $k$) des quantités invariantes en changeant $G$ en $G_{\mathrm{réd}}^\circ$, et de plus pour $\rho_{aff}$ et $d$, c'est une fonction \uline{\unsure{croissante|croissant}} du groupe. Donc cela nous permet [quitte à remplacer $G$ par $\overline{G_1^{\unsure{\circ|a}}}$] de supposer \uline{en tous cas} $G$ plat sur $S$, \deleted{et $G_1$ \unsure{connexe}} et de plus
% archive sheet number: 31
\scanpage{33}
(quitte à faire une extension finie sur $S$ et remplacer $G$ par $\overline{G_{\deleted{\unsure{\mathrm{réd}}}1\,\mathrm{réd}}^{\circ}}$] que $G_1$ est \inserted{connexe et} lisse $/k(s_1)$, enfin [quitte à remplacer $G$ par $G$ $-$ composantes non neutres de $G_0$] que $G_0$ est également connexe [mais peut-être pas lisse].

a) L'inégalité $d\geq d'$ bien connue, mise pour mémoire, elle se réduit \deleted{\unsure{alors}}\inserted{ici} (cas plat) à l'égalité $d=d'$. L'inégalité $\rho_{aff}\geq\rho'_{aff}$ \uline{équivaut alors} à $\rho_{ab}\leq\rho'_{ab}$ [\unsure{car|or} $\rho_{aff}+\rho_{ab}=\rho'_{aff}+\rho'_{ab}=d$], que nous allons traiter plus bas.

b) \deleted{Considérons} On peut supposer que $G_0$ admet un tore maximal splitté $T_0$ (car on peut supposer $k(s_0)$ alg. clos]. \deleted{\unsure{les|Les} $ {}_nT_0$ Soit} $C_0$ le Cartan associé. \deleted{\unsure{si|on peut} \unsure{premier|prolonge|prolonger} \unsure{premier} à \unsure{la caract.|la car.} résiduelle.} Pour tt $n$ premier

\margin{\deleted{$G$ lisse$/S$}}
\deleted{Alors pour les puiss. \unsure{$n\ell$|$\ell$|de $\ell$} de $n$,}
à la caractéristique résiduelle, considérons

\deleted{$ {}_n C_0=\operatorname{Ker}$ [\unsure{mult$^n$|$n^{\mathrm{ième}}$} de $C_0$]. \unsure{Comme} $C_0/T_0=\deleted{\unsure{G}}M_0$ est une extension d'une VA $A_0$ par un groupe unipotent \inserted{lisse} connexe $U_0$, on voit que les extensions \unsure{correspondantes} de $ {}_n A_0$ par $U_0$ splittent \uuline{canoniquement}, \deleted{\unsure{d'où|donc} \unsure{son extension|des extensions}}\inserted{l'image étant nlle, d'où}}
% archive sheet number: 32
\scanpage{34}
$ {}_nT_0$, on peut le prolonger en un sous-groupe \inserted{fini} étale $H(n)$ de $G$, tel que $H(m)={}_{\unsure{m|n}}H(n)$ si $n\mid m$. \margin{$G$ lisse$/k$}La famille des sous-groupes fermés lisses $\underline{\mathrm{Centr}}_G(H(n))$ de $G$ est stationnaire, sa valeur limite $\deleted{D}C$ est telle que $D_0=\mathrm{Centr}_{G_0}(T_0)=C_0$. De plus, on voit que $C_1$ a $\overline{\mathrm m}$ rang réductif et nilpotent que $G_1$ [i.e. \deleted{\unsure{contient}} $C_{\bar{s}_1}^{\deleted{\unsure{\circ}}}$ contient un Cartan de $G_{\bar{s}_1}$] : c'est un \uline{lemme à rapporter}.

\margin{$\ast$}
Il s'ensuit que \deleted{$\rho_n^{\unsure{1}}(G_1)=\rho_n(\unsure{C_1}),\ \rho_r(G_1)=\rho_r(C_1)$} $\rho_r,\rho_u,\rho_{ab},\rho_n,\rho_{na}$ sont les $\overline{\mathrm m}\mathrm{s}$ pour $G_1$ et $C_1$
\[
\begin{aligned}
\deleted{\rho_{\unsure{u|n}}(G_1)=\rho_n(C_1)}\\
\deleted{\rho_u(G_1)=\rho_u(C_1)}\\
\deleted{\rho_{ab}(G_1)=\rho_{ab}(C_1)}\\
\deleted{\rho_r(G_1)=\rho_r(C_1)}
\end{aligned}
\]
et sont aussi les $\overline{\mathrm m}\mathrm{s}$ pour $C_0$ et $G_0$, donc pour les formules 1) à 5), on est ramené au cas où $G=C$, i.e. $G_0$ \deleted{à fibres} est nilpotent. Alors la relation $\boxed{\rho_n\geq\rho'_n}$ \unsure{triviale}, car $\rho_n=d=d'$

\uline{Lemme} $C$ gpe nilpotent connexe lisse sur corps alg. clos $k$. \margin{\unsure{peut|pour} s'énoncer sans alg. clos\ldots}
% slanted
Alors $C/\overbrace{\underline{\mathrm{Centre}}(C)_{\mathrm{réd}}^\circ}^{\unsure{Z=|Z}}$ est unipotent. Pour tt gpe de type multiplicatif $H$, $\underline{\mathrm{Hom}}_{k\text{-gr}}(H,C)\simeq\underline{\mathrm{Hom}}_{k\text{-gr}}(H,Z)$. En particulier, pour $n$ premier à la car, $ {}_nC\simeq{}_nZ$ est un \unsure{groupe|sous-groupe} comm. étale de $C$, central,
% archive sheet number: 33
\scanpage{35}
\deleted{\uline{Corollaire} Soit $G$ lisse sur $S$, et supposons que}

\deleted{On sait qu'un tel ss-gpe \unsure{de|est} $S$}

Appliquons ceci au cas de $C=G_0$; alors, si $S$ est \deleted{\unsure{loc}} str. local, les $ {}_nC$ se prolongent de \uline{façon unique} en des sous-groupes $H(n)$ \deleted{de} étales finis de $G$, \unsure{nécessairement|nécessairemt} centraux, et
\[
{}_mH(n)=H(m)\qquad\text{si }m\mid n
\]

Noter que $\rho_{ab}$ et $\rho_r$, donc $\rho_{ab}+\rho_r=\rho_{\mathrm{rig}}$ sont fonctions croissantes d'un groupe \deleted{lisse} alg.

Soit $K_1\deleted{\unsure{1}}$ l'adhérence de la \unsure{suite|réun$^{\rm ion}$|limite} des $H(n)_1$, \inserted{[qu'on peut supposer lisse$/k(s_1)$ !]} et $\unsure{K|K_1|\overline{K}_1}=\overline{K_1}^{\circ}$, alors \deleted{\unsure{Pour} \unsure{avoir} \unsure{le} \unsure{traiter}}.

Alors
\[
\begin{aligned}
\rho_{ab}(K_1)&\leq\rho_{ab}(G_1)=\rho'_{ab}\\
\rho_{\mathrm{rig}}(K_1)&\leq\rho_{\mathrm{rig}}(\deleted{\unsure{K}}G_1)=\rho'_{\mathrm{rig}}.
\end{aligned}
\]
\inserted{D'autre part, on voit facilement \inserted{grâce à $K_0\supset H(n)_0={}_nG_0$} (\unsure{lemme|comme} ....),}
\[
\begin{aligned}
\inserted{\rho_{ab}(G_0)=\rho_{ab}(K_0)}\\
\inserted{\rho_r(G_0)=\rho_r(K_0)}
\end{aligned}
\]
Donc pour prouver (1) et (2), on est ramené à prouver la $\overline{\mathrm m}$ chose pour $K$ remplaçant $G$.

On a perdu la lissité de $G$, mais on a gagné la \uline{commutativité}, qui implique d'ailleurs $\rho_n=\deleted{\unsure{\rho_n}}\rho'_n=d$. L'inégalité (2) équivaut alors à l'inégalité (6), \deleted{donc on} \inserted{et}
\margin{$\text{\textcircled{\scriptsize 1}}\ \text{à}\ \text{\textcircled{\scriptsize 7}}\ \text{\unsure{n.|ou|et}}\ \text{\textcircled{\scriptsize 8}}$}
% archive sheet number: 34
\scanpage{36}
est ramené à prouver $\unsure{\rho_u|\rho_n}\geq\rho'_u$, \inserted{puis $\deleted{\unsure{\rho}}d_{aff}\geq d'_{aff}$} Remplaçant $G$ par l'adhérence \deleted{de} $G_{1\,\mathrm{unipot}.}$ [qu'on peut supposer défini $/k(s_1)$] on est ramené \inserted{pour $\rho_u\geq\rho'_u$} à prouver que si \uline{$G_1$ est unipotent}, \uline{de $\overline{\mathrm m}$} $G_0$.
\margin{[sans restriction de lissité]}
% slanted

Facile par les points d'ordre fini.

Pour la deuxième assertion, on est ramené à prouver que \uline{si} $G_1$ \uline{est affine}, $G_0$ \uline{est affine}.
\margin{[sans restriction de lissité]}
% slanted

On est ramené au cas \deleted{\unsure{d'une}} où $S$ est \deleted{\unsure{un}} schéma \inserted{intègre} de \uline{type fini} \deleted{sur} \deleted{\unsure{le}} $\mathbb Z$, \inserted{$S$ \unsure{normal|ouvert|semi-local}, $\dim S=1$} (à préciser \ldots).

\deleted{Considérons les $ {}_nG$ et la partie \unsure{compos.} première.} On peut supposer \deleted{\unsure{le groupe}} \uline{le tore maximal de $G_1$ est splitté}. $s_0$ pt fermé, $s_1$ pt génériq. Donc en tous cas $k(s_0)$ est un corps fini.

\diagram{\begin{tikzcd}[ampersand replacement=\&,row sep=small]
G\arrow[d]\\S
\end{tikzcd}}
\figure{\begin{tikzpicture}[x=1cm,y=1cm]
\draw (0.9,1.5) .. controls (1.3,1.43) and (1.65,1.55) .. (1.85,1.51);
\draw[->] (1.76,1.49) .. controls (1.85,1.49) and (1.86,1.50) .. (1.89,1.46);
\draw[->] (1.85,1.04) -- (0.68,1.04) -- (0.68,.62);
\fill (.99,.94) circle (.035);
\draw[thick] (-.47,-1.16) .. controls (-.10,-1.04) and (.28,-.92) .. (.85,-1.10);
\fill (.13,-.99) circle (.045);
\end{tikzpicture}}

Passant à la cloture hensélienne (mais pas strictement hensélienne !) en $s_0$, d'où un $\widetilde S$, $\widetilde G$ sur $\widetilde S$, \deleted{les $ {}_n\widetilde G$} Soit $l$ distinct de la caractéristique résiduelle, $n$ une puissance variable de $l$, $ {}_n\widetilde G$ a une partie \unsure{\uline{finie}|$\widetilde C^{\rm rie}$} sur $\widetilde S$, soit $\widetilde H(n)$.

Alors $\pi_1(\widetilde S)=\pi_1(s_0)\simeq\mathbb Z$ opère sur $H(n)_{\bar s_0}=M_n$, d'où par passage à la limite des opérations de $\mathbb Z$ sur $\varprojlim M_n=M\simeq\mathbb Z_l^{2\rho_{ab}+\rho_r}$. D'une part, \deleted{ces opérations donnent} comme $M_\infty\simeq\varprojlim(\deleted{\unsure{{}_nG}}\inserted{H(n)})_{\bar s_1}$ \deleted{\unsure{ce}} [\unsure{l'immersion|immersion|opération} détermine mod une opération de $\pi_1(\widetilde S)$ sur $M_\infty$] d'où $M_\infty\subset\varprojlim_n({}_nG)_{\bar s_1}=$ \deleted{donc}
% archive sheet number: 35
\scanpage{37}
$\varprojlim_n({}_nT)_{\bar s_1}=({}_nG_m)^r_{\bar s_1}$, donc les opérations de $\pi_1(\widetilde S,\bar s_0)$ sur $M_\infty$ se font par l'intermédiaire du ``caractère \unsure{standart|standard}'' l'opération sur les racines de l'unité. Or si on a $\rho_{ab}\neq 0$, alors les opérations du gpe de Galois $\pi_{\deleted{\unsure{0}}1}(\widetilde S,\bar s_0)=\pi_{\unsure{1|0}}(s_0,\bar s_0)$ \uline{ne se font pas} par l'intermédiaire exclusif de ce caractère \ldots

On est ramené à un énoncé sur une V.A. sur un corps \deleted{\unsure{f|$\mathfrak{k}$}} de type fini [fini serait ici suffisant \ldots]

\begin{center}\rule{2.5cm}{0.35pt}\end{center}

Restent à prouver $\text{\textcircled{\scriptsize 3}}$ et $\text{\textcircled{\scriptsize\unsure{4|9}}}$.
\margin{\unsure{commencer} par cela ....}
% slanted

$\text{\textcircled{\scriptsize 3}}\Longleftrightarrow\text{\textcircled{\scriptsize 9}}$, et \deleted{\unsure{résulte}} \unsure{résulterait|résulte|résultat} du fait qu'\uline{un gpe résoluble se spécialise en un résoluble}

$\text{\textcircled{\scriptsize 4}}$ On regarde un tore maximal $T_0$ de $G_0$
\diagram{\begin{tikzcd}[ampersand replacement=\&,row sep=0pt,column sep=1em]
G\&G_1\\
\cup\&\cup\\
\overline{R_1}=R_0\&R_1
\end{tikzcd}}
\margin{$G_0/R_0$}
on prolonge les $ {}_nT$ en des $H(n)$ de façon compatible, alors on voit \unsure{(si|si|$G_1$} $r_0=\rho_{ss}(G_0)\inserted{\leq\dim T_0/T_0\cap R_0}$) que les $H(n)_1/R_1\cap H(n)_1$ sont de ``$n$-rang'' $\geq\dim T_0/T_0\cap R_0\geq r_0$. Ils sont de $n$-rang $\leq\rho'_{ss}$, d'où $\rho'_{ss}\geq\rho_{ss}$ \unsure{\textendash|\ldots}
% archive sheet number: 36
